\documentclass{beamer}

\usepackage[utf8]{inputenc}
\input{../helper}

\title{Introdução à Computação \\ tipos abstratos, filas}
\author{Alexandre Rademaker}
\date{}

\begin{document}

\begin{frame}
 \maketitle
\end{frame}


\begin{frame}[allowframebreaks,fragile]{filas}

  Outro tipo abstrato de dados é a \textbf{fila}, onde valores são
  colocados em uma extremidade e retirados da extremidade oposta.

  \begin{center}
    \includegraphics[width=.6\textwidth]{queue.png}
  \end{center}

  \framebreak

  Fila um tipo abstrato de dados do tipo ``First-In, First Out''
  (FIFO), ou seja ``o primeiro que entra é o primeiro que sai''.

  Filas podem ser implementadas com listas encadeadas, arrays, arrays
  circulares. 

  \framebreak

  Filas são usadas em várias áreas da computação onde o acesso a algum
  recurso é controlado e requisições devem aguardar para serem
  atendidas: CPU, fila de impressão, roteadores etc.

  Na navegação em grafos, \emph{buscas em profundidade} são
  implementadas com \emph{pilhas} e \emph{buscas em largura} com
  \emph{filas}.

  \framebreak

  As duas principais operações com filas são:

  \begin{description}
  \item[dequeue] retira um elemento da frente da fila. Na
    implementação com arrays pode ser $O(n)$ mas com listas encadeadas
    é $O(1)$.
  \item[enqueue] insere um elemento no final da fila. Complexidade
    $O(1)$.
  \end{description}

  Outras operações possíveis (e auxiliares): imprimir a fila,
  inspecionar o primeiro elemento, inspecionar o último elemento,
  determinar se a fila está vazia, tamanho da fila, etc.

  \framebreak

  Em C, vamos implementar uma \texttt{struct} contendo dois ponteiros,
  um para a \emph{frente} e outro para o \emph{final} da fila.

  \begin{center}
    \includegraphics[width=.9\textwidth]{queue-c.png}
  \end{center}

\end{frame}


\begin{frame}[allowframebreaks,fragile]{filas em C}
  
  Além da estrutura já criada para listas encadeadas, uma estrutura
  adicional para conter os o ponteiro para o início da fila (front) e
  final da fila (rear).

  \begin{lstlisting}[style=normalc,caption=src/queue-list.c]
    typedef struct node {
      int number;
      struct node *next;
    } t_node;
    
    typedef struct queue {
      t_node *front;
      t_node *rear;
    } t_queue;
  \end{lstlisting}

  \framebreak

  \begin{lstlisting}[style=normalc,caption=src/queue-list.c]
    void enqueue(t_queue* q, int k) {
      t_node* temp = create_node(k);

      // if empty
      if (q->rear == NULL) {
        q->front = q->rear = temp;
        return;
      }
      
      q->rear->next = temp;
      q->rear = temp;
    }
  \end{lstlisting}

  \framebreak

  \begin{lstlisting}[style=normalc,caption=src/queue-list.c]
    void dequeue(t_queue* q) {
      // if empty
      if (q->front == NULL) return;

      // save previous front and move
      t_node* temp = q->front;
      q->front = q->front->next;

      // if front becomes NULL
      if (q->front == NULL)
        q->rear = NULL;
      
      free(temp);
    }
  \end{lstlisting}

  \framebreak

  Using it

  \begin{lstlisting}[style=normalc]
    t_queue* create_queue() {
      t_queue *q = malloc(sizeof(t_queue));
      q->front = q->rear = NULL;
      return q;
    }
    
    int main(int argc, char *argv[]) {
      t_queue *q = create_queue();
      enqueue(q, 1);
      enqueue(q, 2);
      dequeue(q);
      enqueue(q, 3);
      enqueue(q, 4);
      dequeue(q);
    }
  \end{lstlisting}
  
\end{frame}


\end{document}

%%% Local Variables:
%%% mode: latex
%%% TeX-master: t
%%% End:
